% spidoc.tex - macros used by the generated reference section.
% Load from the preamble with % spidoc.tex - macros used by the generated reference section.
% Load from the preamble with % spidoc.tex - macros used by the generated reference section.
% Load from the preamble with % spidoc.tex - macros used by the generated reference section.
% Load from the preamble with \input{spidoc}.
% Needs the longtable and fancyvrb packages, which the main document loads.

% General table.
%   #1 number of columns, #2 column specification, #3 heading row
%
% The table has one extra empty column on the right, and the rules are drawn
% with \cline over the real columns only. Rows written by hand may therefore
% end with "& \\" as they always did. (Once those rows have been tidied up the
% extra column can be dropped and \cline replaced by \hline.)
\newcounter{spicols}
\newenvironment{spitable}[3]{%
  \setcounter{spicols}{#1}%
  \small
  \begin{longtable}{#2l}
  \cline{1-#1} #3 \\ \endfirsthead
  \cline{1-#1} #3 \\ \cline{1-#1} \endhead
  \cline{1-#1} \multicolumn{#1}{r}{\emph{Continued on next page}} \\ \endfoot
  \cline{1-#1} \endlastfoot
}{%
  \end{longtable}%
}

% One row of a general table: the cells separated by &.
\newcommand{\spirow}[1]{\cline{1-\value{spicols}} #1 \\}

% Options table: Name / Default / Range.
\newenvironment{optionstable}{%
  \begin{spitable}{3}{|p{80pt}|p{150pt}|p{140pt}|}%
    {\textbf{Name} & \textbf{Default} & \textbf{Range}}%
}{%
  \end{spitable}%
}
\newcommand{\optrow}[3]{\spirow{\texttt{#1} & #2 & #3}}

% Syntax block: verbatim, but \ { } still work as commands inside it.
% This needs the fancyvrb package; documents which only use tables do not
% need to load it, so the environment is only defined when it is available.
\ifdefined\Verbatim
\DefineVerbatimEnvironment{syntax}{Verbatim}{commandchars=\\\{\}}
\fi
\newcommand{\syntaxhead}{\par\textbf{SYNTAX:}\nopagebreak}
\newcommand{\syntaxdesc}{\emph{/** description */}}

% A variable name such as $returnType.
\newcommand{\spival}[1]{\texttt{\$#1}}

% Marks the macros as loaded, so generated files can load this file on demand:
%   \ifdefined\spidocloaded\else\input{spidoc}\fi
\newcommand{\spidocloaded}{}

.
% Needs the longtable and fancyvrb packages, which the main document loads.

% General table.
%   #1 number of columns, #2 column specification, #3 heading row
%
% The table has one extra empty column on the right, and the rules are drawn
% with \cline over the real columns only. Rows written by hand may therefore
% end with "& \\" as they always did. (Once those rows have been tidied up the
% extra column can be dropped and \cline replaced by \hline.)
\newcounter{spicols}
\newenvironment{spitable}[3]{%
  \setcounter{spicols}{#1}%
  \small
  \begin{longtable}{#2l}
  \cline{1-#1} #3 \\ \endfirsthead
  \cline{1-#1} #3 \\ \cline{1-#1} \endhead
  \cline{1-#1} \multicolumn{#1}{r}{\emph{Continued on next page}} \\ \endfoot
  \cline{1-#1} \endlastfoot
}{%
  \end{longtable}%
}

% One row of a general table: the cells separated by &.
\newcommand{\spirow}[1]{\cline{1-\value{spicols}} #1 \\}

% Options table: Name / Default / Range.
\newenvironment{optionstable}{%
  \begin{spitable}{3}{|p{80pt}|p{150pt}|p{140pt}|}%
    {\textbf{Name} & \textbf{Default} & \textbf{Range}}%
}{%
  \end{spitable}%
}
\newcommand{\optrow}[3]{\spirow{\texttt{#1} & #2 & #3}}

% Syntax block: verbatim, but \ { } still work as commands inside it.
% This needs the fancyvrb package; documents which only use tables do not
% need to load it, so the environment is only defined when it is available.
\ifdefined\Verbatim
\DefineVerbatimEnvironment{syntax}{Verbatim}{commandchars=\\\{\}}
\fi
\newcommand{\syntaxhead}{\par\textbf{SYNTAX:}\nopagebreak}
\newcommand{\syntaxdesc}{\emph{/** description */}}

% A variable name such as $returnType.
\newcommand{\spival}[1]{\texttt{\$#1}}

% Marks the macros as loaded, so generated files can load this file on demand:
%   \ifdefined\spidocloaded\else% spidoc.tex - macros used by the generated reference section.
% Load from the preamble with \input{spidoc}.
% Needs the longtable and fancyvrb packages, which the main document loads.

% General table.
%   #1 number of columns, #2 column specification, #3 heading row
%
% The table has one extra empty column on the right, and the rules are drawn
% with \cline over the real columns only. Rows written by hand may therefore
% end with "& \\" as they always did. (Once those rows have been tidied up the
% extra column can be dropped and \cline replaced by \hline.)
\newcounter{spicols}
\newenvironment{spitable}[3]{%
  \setcounter{spicols}{#1}%
  \small
  \begin{longtable}{#2l}
  \cline{1-#1} #3 \\ \endfirsthead
  \cline{1-#1} #3 \\ \cline{1-#1} \endhead
  \cline{1-#1} \multicolumn{#1}{r}{\emph{Continued on next page}} \\ \endfoot
  \cline{1-#1} \endlastfoot
}{%
  \end{longtable}%
}

% One row of a general table: the cells separated by &.
\newcommand{\spirow}[1]{\cline{1-\value{spicols}} #1 \\}

% Options table: Name / Default / Range.
\newenvironment{optionstable}{%
  \begin{spitable}{3}{|p{80pt}|p{150pt}|p{140pt}|}%
    {\textbf{Name} & \textbf{Default} & \textbf{Range}}%
}{%
  \end{spitable}%
}
\newcommand{\optrow}[3]{\spirow{\texttt{#1} & #2 & #3}}

% Syntax block: verbatim, but \ { } still work as commands inside it.
% This needs the fancyvrb package; documents which only use tables do not
% need to load it, so the environment is only defined when it is available.
\ifdefined\Verbatim
\DefineVerbatimEnvironment{syntax}{Verbatim}{commandchars=\\\{\}}
\fi
\newcommand{\syntaxhead}{\par\textbf{SYNTAX:}\nopagebreak}
\newcommand{\syntaxdesc}{\emph{/** description */}}

% A variable name such as $returnType.
\newcommand{\spival}[1]{\texttt{\$#1}}

% Marks the macros as loaded, so generated files can load this file on demand:
%   \ifdefined\spidocloaded\else\input{spidoc}\fi
\newcommand{\spidocloaded}{}

\fi
\newcommand{\spidocloaded}{}

.
% Needs the longtable and fancyvrb packages, which the main document loads.

% General table.
%   #1 number of columns, #2 column specification, #3 heading row
%
% The table has one extra empty column on the right, and the rules are drawn
% with \cline over the real columns only. Rows written by hand may therefore
% end with "& \\" as they always did. (Once those rows have been tidied up the
% extra column can be dropped and \cline replaced by \hline.)
\newcounter{spicols}
\newenvironment{spitable}[3]{%
  \setcounter{spicols}{#1}%
  \small
  \begin{longtable}{#2l}
  \cline{1-#1} #3 \\ \endfirsthead
  \cline{1-#1} #3 \\ \cline{1-#1} \endhead
  \cline{1-#1} \multicolumn{#1}{r}{\emph{Continued on next page}} \\ \endfoot
  \cline{1-#1} \endlastfoot
}{%
  \end{longtable}%
}

% One row of a general table: the cells separated by &.
\newcommand{\spirow}[1]{\cline{1-\value{spicols}} #1 \\}

% Options table: Name / Default / Range.
\newenvironment{optionstable}{%
  \begin{spitable}{3}{|p{80pt}|p{150pt}|p{140pt}|}%
    {\textbf{Name} & \textbf{Default} & \textbf{Range}}%
}{%
  \end{spitable}%
}
\newcommand{\optrow}[3]{\spirow{\texttt{#1} & #2 & #3}}

% Syntax block: verbatim, but \ { } still work as commands inside it.
% This needs the fancyvrb package; documents which only use tables do not
% need to load it, so the environment is only defined when it is available.
\ifdefined\Verbatim
\DefineVerbatimEnvironment{syntax}{Verbatim}{commandchars=\\\{\}}
\fi
\newcommand{\syntaxhead}{\par\textbf{SYNTAX:}\nopagebreak}
\newcommand{\syntaxdesc}{\emph{/** description */}}

% A variable name such as $returnType.
\newcommand{\spival}[1]{\texttt{\$#1}}

% Marks the macros as loaded, so generated files can load this file on demand:
%   \ifdefined\spidocloaded\else% spidoc.tex - macros used by the generated reference section.
% Load from the preamble with % spidoc.tex - macros used by the generated reference section.
% Load from the preamble with \input{spidoc}.
% Needs the longtable and fancyvrb packages, which the main document loads.

% General table.
%   #1 number of columns, #2 column specification, #3 heading row
%
% The table has one extra empty column on the right, and the rules are drawn
% with \cline over the real columns only. Rows written by hand may therefore
% end with "& \\" as they always did. (Once those rows have been tidied up the
% extra column can be dropped and \cline replaced by \hline.)
\newcounter{spicols}
\newenvironment{spitable}[3]{%
  \setcounter{spicols}{#1}%
  \small
  \begin{longtable}{#2l}
  \cline{1-#1} #3 \\ \endfirsthead
  \cline{1-#1} #3 \\ \cline{1-#1} \endhead
  \cline{1-#1} \multicolumn{#1}{r}{\emph{Continued on next page}} \\ \endfoot
  \cline{1-#1} \endlastfoot
}{%
  \end{longtable}%
}

% One row of a general table: the cells separated by &.
\newcommand{\spirow}[1]{\cline{1-\value{spicols}} #1 \\}

% Options table: Name / Default / Range.
\newenvironment{optionstable}{%
  \begin{spitable}{3}{|p{80pt}|p{150pt}|p{140pt}|}%
    {\textbf{Name} & \textbf{Default} & \textbf{Range}}%
}{%
  \end{spitable}%
}
\newcommand{\optrow}[3]{\spirow{\texttt{#1} & #2 & #3}}

% Syntax block: verbatim, but \ { } still work as commands inside it.
% This needs the fancyvrb package; documents which only use tables do not
% need to load it, so the environment is only defined when it is available.
\ifdefined\Verbatim
\DefineVerbatimEnvironment{syntax}{Verbatim}{commandchars=\\\{\}}
\fi
\newcommand{\syntaxhead}{\par\textbf{SYNTAX:}\nopagebreak}
\newcommand{\syntaxdesc}{\emph{/** description */}}

% A variable name such as $returnType.
\newcommand{\spival}[1]{\texttt{\$#1}}

% Marks the macros as loaded, so generated files can load this file on demand:
%   \ifdefined\spidocloaded\else\input{spidoc}\fi
\newcommand{\spidocloaded}{}

.
% Needs the longtable and fancyvrb packages, which the main document loads.

% General table.
%   #1 number of columns, #2 column specification, #3 heading row
%
% The table has one extra empty column on the right, and the rules are drawn
% with \cline over the real columns only. Rows written by hand may therefore
% end with "& \\" as they always did. (Once those rows have been tidied up the
% extra column can be dropped and \cline replaced by \hline.)
\newcounter{spicols}
\newenvironment{spitable}[3]{%
  \setcounter{spicols}{#1}%
  \small
  \begin{longtable}{#2l}
  \cline{1-#1} #3 \\ \endfirsthead
  \cline{1-#1} #3 \\ \cline{1-#1} \endhead
  \cline{1-#1} \multicolumn{#1}{r}{\emph{Continued on next page}} \\ \endfoot
  \cline{1-#1} \endlastfoot
}{%
  \end{longtable}%
}

% One row of a general table: the cells separated by &.
\newcommand{\spirow}[1]{\cline{1-\value{spicols}} #1 \\}

% Options table: Name / Default / Range.
\newenvironment{optionstable}{%
  \begin{spitable}{3}{|p{80pt}|p{150pt}|p{140pt}|}%
    {\textbf{Name} & \textbf{Default} & \textbf{Range}}%
}{%
  \end{spitable}%
}
\newcommand{\optrow}[3]{\spirow{\texttt{#1} & #2 & #3}}

% Syntax block: verbatim, but \ { } still work as commands inside it.
% This needs the fancyvrb package; documents which only use tables do not
% need to load it, so the environment is only defined when it is available.
\ifdefined\Verbatim
\DefineVerbatimEnvironment{syntax}{Verbatim}{commandchars=\\\{\}}
\fi
\newcommand{\syntaxhead}{\par\textbf{SYNTAX:}\nopagebreak}
\newcommand{\syntaxdesc}{\emph{/** description */}}

% A variable name such as $returnType.
\newcommand{\spival}[1]{\texttt{\$#1}}

% Marks the macros as loaded, so generated files can load this file on demand:
%   \ifdefined\spidocloaded\else% spidoc.tex - macros used by the generated reference section.
% Load from the preamble with \input{spidoc}.
% Needs the longtable and fancyvrb packages, which the main document loads.

% General table.
%   #1 number of columns, #2 column specification, #3 heading row
%
% The table has one extra empty column on the right, and the rules are drawn
% with \cline over the real columns only. Rows written by hand may therefore
% end with "& \\" as they always did. (Once those rows have been tidied up the
% extra column can be dropped and \cline replaced by \hline.)
\newcounter{spicols}
\newenvironment{spitable}[3]{%
  \setcounter{spicols}{#1}%
  \small
  \begin{longtable}{#2l}
  \cline{1-#1} #3 \\ \endfirsthead
  \cline{1-#1} #3 \\ \cline{1-#1} \endhead
  \cline{1-#1} \multicolumn{#1}{r}{\emph{Continued on next page}} \\ \endfoot
  \cline{1-#1} \endlastfoot
}{%
  \end{longtable}%
}

% One row of a general table: the cells separated by &.
\newcommand{\spirow}[1]{\cline{1-\value{spicols}} #1 \\}

% Options table: Name / Default / Range.
\newenvironment{optionstable}{%
  \begin{spitable}{3}{|p{80pt}|p{150pt}|p{140pt}|}%
    {\textbf{Name} & \textbf{Default} & \textbf{Range}}%
}{%
  \end{spitable}%
}
\newcommand{\optrow}[3]{\spirow{\texttt{#1} & #2 & #3}}

% Syntax block: verbatim, but \ { } still work as commands inside it.
% This needs the fancyvrb package; documents which only use tables do not
% need to load it, so the environment is only defined when it is available.
\ifdefined\Verbatim
\DefineVerbatimEnvironment{syntax}{Verbatim}{commandchars=\\\{\}}
\fi
\newcommand{\syntaxhead}{\par\textbf{SYNTAX:}\nopagebreak}
\newcommand{\syntaxdesc}{\emph{/** description */}}

% A variable name such as $returnType.
\newcommand{\spival}[1]{\texttt{\$#1}}

% Marks the macros as loaded, so generated files can load this file on demand:
%   \ifdefined\spidocloaded\else\input{spidoc}\fi
\newcommand{\spidocloaded}{}

\fi
\newcommand{\spidocloaded}{}

\fi
\newcommand{\spidocloaded}{}

.
% Needs the longtable and fancyvrb packages, which the main document loads.

% General table.
%   #1 number of columns, #2 column specification, #3 heading row
%
% The table has one extra empty column on the right, and the rules are drawn
% with \cline over the real columns only. Rows written by hand may therefore
% end with "& \\" as they always did. (Once those rows have been tidied up the
% extra column can be dropped and \cline replaced by \hline.)
\newcounter{spicols}
\newenvironment{spitable}[3]{%
  \setcounter{spicols}{#1}%
  \small
  \begin{longtable}{#2l}
  \cline{1-#1} #3 \\ \endfirsthead
  \cline{1-#1} #3 \\ \cline{1-#1} \endhead
  \cline{1-#1} \multicolumn{#1}{r}{\emph{Continued on next page}} \\ \endfoot
  \cline{1-#1} \endlastfoot
}{%
  \end{longtable}%
}

% One row of a general table: the cells separated by &.
\newcommand{\spirow}[1]{\cline{1-\value{spicols}} #1 \\}

% Options table: Name / Default / Range.
\newenvironment{optionstable}{%
  \begin{spitable}{3}{|p{80pt}|p{150pt}|p{140pt}|}%
    {\textbf{Name} & \textbf{Default} & \textbf{Range}}%
}{%
  \end{spitable}%
}
\newcommand{\optrow}[3]{\spirow{\texttt{#1} & #2 & #3}}

% Syntax block: verbatim, but \ { } still work as commands inside it.
% This needs the fancyvrb package; documents which only use tables do not
% need to load it, so the environment is only defined when it is available.
\ifdefined\Verbatim
\DefineVerbatimEnvironment{syntax}{Verbatim}{commandchars=\\\{\}}
\fi
\newcommand{\syntaxhead}{\par\textbf{SYNTAX:}\nopagebreak}
\newcommand{\syntaxdesc}{\emph{/** description */}}

% A variable name such as $returnType.
\newcommand{\spival}[1]{\texttt{\$#1}}

% Marks the macros as loaded, so generated files can load this file on demand:
%   \ifdefined\spidocloaded\else% spidoc.tex - macros used by the generated reference section.
% Load from the preamble with % spidoc.tex - macros used by the generated reference section.
% Load from the preamble with % spidoc.tex - macros used by the generated reference section.
% Load from the preamble with \input{spidoc}.
% Needs the longtable and fancyvrb packages, which the main document loads.

% General table.
%   #1 number of columns, #2 column specification, #3 heading row
%
% The table has one extra empty column on the right, and the rules are drawn
% with \cline over the real columns only. Rows written by hand may therefore
% end with "& \\" as they always did. (Once those rows have been tidied up the
% extra column can be dropped and \cline replaced by \hline.)
\newcounter{spicols}
\newenvironment{spitable}[3]{%
  \setcounter{spicols}{#1}%
  \small
  \begin{longtable}{#2l}
  \cline{1-#1} #3 \\ \endfirsthead
  \cline{1-#1} #3 \\ \cline{1-#1} \endhead
  \cline{1-#1} \multicolumn{#1}{r}{\emph{Continued on next page}} \\ \endfoot
  \cline{1-#1} \endlastfoot
}{%
  \end{longtable}%
}

% One row of a general table: the cells separated by &.
\newcommand{\spirow}[1]{\cline{1-\value{spicols}} #1 \\}

% Options table: Name / Default / Range.
\newenvironment{optionstable}{%
  \begin{spitable}{3}{|p{80pt}|p{150pt}|p{140pt}|}%
    {\textbf{Name} & \textbf{Default} & \textbf{Range}}%
}{%
  \end{spitable}%
}
\newcommand{\optrow}[3]{\spirow{\texttt{#1} & #2 & #3}}

% Syntax block: verbatim, but \ { } still work as commands inside it.
% This needs the fancyvrb package; documents which only use tables do not
% need to load it, so the environment is only defined when it is available.
\ifdefined\Verbatim
\DefineVerbatimEnvironment{syntax}{Verbatim}{commandchars=\\\{\}}
\fi
\newcommand{\syntaxhead}{\par\textbf{SYNTAX:}\nopagebreak}
\newcommand{\syntaxdesc}{\emph{/** description */}}

% A variable name such as $returnType.
\newcommand{\spival}[1]{\texttt{\$#1}}

% Marks the macros as loaded, so generated files can load this file on demand:
%   \ifdefined\spidocloaded\else\input{spidoc}\fi
\newcommand{\spidocloaded}{}

.
% Needs the longtable and fancyvrb packages, which the main document loads.

% General table.
%   #1 number of columns, #2 column specification, #3 heading row
%
% The table has one extra empty column on the right, and the rules are drawn
% with \cline over the real columns only. Rows written by hand may therefore
% end with "& \\" as they always did. (Once those rows have been tidied up the
% extra column can be dropped and \cline replaced by \hline.)
\newcounter{spicols}
\newenvironment{spitable}[3]{%
  \setcounter{spicols}{#1}%
  \small
  \begin{longtable}{#2l}
  \cline{1-#1} #3 \\ \endfirsthead
  \cline{1-#1} #3 \\ \cline{1-#1} \endhead
  \cline{1-#1} \multicolumn{#1}{r}{\emph{Continued on next page}} \\ \endfoot
  \cline{1-#1} \endlastfoot
}{%
  \end{longtable}%
}

% One row of a general table: the cells separated by &.
\newcommand{\spirow}[1]{\cline{1-\value{spicols}} #1 \\}

% Options table: Name / Default / Range.
\newenvironment{optionstable}{%
  \begin{spitable}{3}{|p{80pt}|p{150pt}|p{140pt}|}%
    {\textbf{Name} & \textbf{Default} & \textbf{Range}}%
}{%
  \end{spitable}%
}
\newcommand{\optrow}[3]{\spirow{\texttt{#1} & #2 & #3}}

% Syntax block: verbatim, but \ { } still work as commands inside it.
% This needs the fancyvrb package; documents which only use tables do not
% need to load it, so the environment is only defined when it is available.
\ifdefined\Verbatim
\DefineVerbatimEnvironment{syntax}{Verbatim}{commandchars=\\\{\}}
\fi
\newcommand{\syntaxhead}{\par\textbf{SYNTAX:}\nopagebreak}
\newcommand{\syntaxdesc}{\emph{/** description */}}

% A variable name such as $returnType.
\newcommand{\spival}[1]{\texttt{\$#1}}

% Marks the macros as loaded, so generated files can load this file on demand:
%   \ifdefined\spidocloaded\else% spidoc.tex - macros used by the generated reference section.
% Load from the preamble with \input{spidoc}.
% Needs the longtable and fancyvrb packages, which the main document loads.

% General table.
%   #1 number of columns, #2 column specification, #3 heading row
%
% The table has one extra empty column on the right, and the rules are drawn
% with \cline over the real columns only. Rows written by hand may therefore
% end with "& \\" as they always did. (Once those rows have been tidied up the
% extra column can be dropped and \cline replaced by \hline.)
\newcounter{spicols}
\newenvironment{spitable}[3]{%
  \setcounter{spicols}{#1}%
  \small
  \begin{longtable}{#2l}
  \cline{1-#1} #3 \\ \endfirsthead
  \cline{1-#1} #3 \\ \cline{1-#1} \endhead
  \cline{1-#1} \multicolumn{#1}{r}{\emph{Continued on next page}} \\ \endfoot
  \cline{1-#1} \endlastfoot
}{%
  \end{longtable}%
}

% One row of a general table: the cells separated by &.
\newcommand{\spirow}[1]{\cline{1-\value{spicols}} #1 \\}

% Options table: Name / Default / Range.
\newenvironment{optionstable}{%
  \begin{spitable}{3}{|p{80pt}|p{150pt}|p{140pt}|}%
    {\textbf{Name} & \textbf{Default} & \textbf{Range}}%
}{%
  \end{spitable}%
}
\newcommand{\optrow}[3]{\spirow{\texttt{#1} & #2 & #3}}

% Syntax block: verbatim, but \ { } still work as commands inside it.
% This needs the fancyvrb package; documents which only use tables do not
% need to load it, so the environment is only defined when it is available.
\ifdefined\Verbatim
\DefineVerbatimEnvironment{syntax}{Verbatim}{commandchars=\\\{\}}
\fi
\newcommand{\syntaxhead}{\par\textbf{SYNTAX:}\nopagebreak}
\newcommand{\syntaxdesc}{\emph{/** description */}}

% A variable name such as $returnType.
\newcommand{\spival}[1]{\texttt{\$#1}}

% Marks the macros as loaded, so generated files can load this file on demand:
%   \ifdefined\spidocloaded\else\input{spidoc}\fi
\newcommand{\spidocloaded}{}

\fi
\newcommand{\spidocloaded}{}

.
% Needs the longtable and fancyvrb packages, which the main document loads.

% General table.
%   #1 number of columns, #2 column specification, #3 heading row
%
% The table has one extra empty column on the right, and the rules are drawn
% with \cline over the real columns only. Rows written by hand may therefore
% end with "& \\" as they always did. (Once those rows have been tidied up the
% extra column can be dropped and \cline replaced by \hline.)
\newcounter{spicols}
\newenvironment{spitable}[3]{%
  \setcounter{spicols}{#1}%
  \small
  \begin{longtable}{#2l}
  \cline{1-#1} #3 \\ \endfirsthead
  \cline{1-#1} #3 \\ \cline{1-#1} \endhead
  \cline{1-#1} \multicolumn{#1}{r}{\emph{Continued on next page}} \\ \endfoot
  \cline{1-#1} \endlastfoot
}{%
  \end{longtable}%
}

% One row of a general table: the cells separated by &.
\newcommand{\spirow}[1]{\cline{1-\value{spicols}} #1 \\}

% Options table: Name / Default / Range.
\newenvironment{optionstable}{%
  \begin{spitable}{3}{|p{80pt}|p{150pt}|p{140pt}|}%
    {\textbf{Name} & \textbf{Default} & \textbf{Range}}%
}{%
  \end{spitable}%
}
\newcommand{\optrow}[3]{\spirow{\texttt{#1} & #2 & #3}}

% Syntax block: verbatim, but \ { } still work as commands inside it.
% This needs the fancyvrb package; documents which only use tables do not
% need to load it, so the environment is only defined when it is available.
\ifdefined\Verbatim
\DefineVerbatimEnvironment{syntax}{Verbatim}{commandchars=\\\{\}}
\fi
\newcommand{\syntaxhead}{\par\textbf{SYNTAX:}\nopagebreak}
\newcommand{\syntaxdesc}{\emph{/** description */}}

% A variable name such as $returnType.
\newcommand{\spival}[1]{\texttt{\$#1}}

% Marks the macros as loaded, so generated files can load this file on demand:
%   \ifdefined\spidocloaded\else% spidoc.tex - macros used by the generated reference section.
% Load from the preamble with % spidoc.tex - macros used by the generated reference section.
% Load from the preamble with \input{spidoc}.
% Needs the longtable and fancyvrb packages, which the main document loads.

% General table.
%   #1 number of columns, #2 column specification, #3 heading row
%
% The table has one extra empty column on the right, and the rules are drawn
% with \cline over the real columns only. Rows written by hand may therefore
% end with "& \\" as they always did. (Once those rows have been tidied up the
% extra column can be dropped and \cline replaced by \hline.)
\newcounter{spicols}
\newenvironment{spitable}[3]{%
  \setcounter{spicols}{#1}%
  \small
  \begin{longtable}{#2l}
  \cline{1-#1} #3 \\ \endfirsthead
  \cline{1-#1} #3 \\ \cline{1-#1} \endhead
  \cline{1-#1} \multicolumn{#1}{r}{\emph{Continued on next page}} \\ \endfoot
  \cline{1-#1} \endlastfoot
}{%
  \end{longtable}%
}

% One row of a general table: the cells separated by &.
\newcommand{\spirow}[1]{\cline{1-\value{spicols}} #1 \\}

% Options table: Name / Default / Range.
\newenvironment{optionstable}{%
  \begin{spitable}{3}{|p{80pt}|p{150pt}|p{140pt}|}%
    {\textbf{Name} & \textbf{Default} & \textbf{Range}}%
}{%
  \end{spitable}%
}
\newcommand{\optrow}[3]{\spirow{\texttt{#1} & #2 & #3}}

% Syntax block: verbatim, but \ { } still work as commands inside it.
% This needs the fancyvrb package; documents which only use tables do not
% need to load it, so the environment is only defined when it is available.
\ifdefined\Verbatim
\DefineVerbatimEnvironment{syntax}{Verbatim}{commandchars=\\\{\}}
\fi
\newcommand{\syntaxhead}{\par\textbf{SYNTAX:}\nopagebreak}
\newcommand{\syntaxdesc}{\emph{/** description */}}

% A variable name such as $returnType.
\newcommand{\spival}[1]{\texttt{\$#1}}

% Marks the macros as loaded, so generated files can load this file on demand:
%   \ifdefined\spidocloaded\else\input{spidoc}\fi
\newcommand{\spidocloaded}{}

.
% Needs the longtable and fancyvrb packages, which the main document loads.

% General table.
%   #1 number of columns, #2 column specification, #3 heading row
%
% The table has one extra empty column on the right, and the rules are drawn
% with \cline over the real columns only. Rows written by hand may therefore
% end with "& \\" as they always did. (Once those rows have been tidied up the
% extra column can be dropped and \cline replaced by \hline.)
\newcounter{spicols}
\newenvironment{spitable}[3]{%
  \setcounter{spicols}{#1}%
  \small
  \begin{longtable}{#2l}
  \cline{1-#1} #3 \\ \endfirsthead
  \cline{1-#1} #3 \\ \cline{1-#1} \endhead
  \cline{1-#1} \multicolumn{#1}{r}{\emph{Continued on next page}} \\ \endfoot
  \cline{1-#1} \endlastfoot
}{%
  \end{longtable}%
}

% One row of a general table: the cells separated by &.
\newcommand{\spirow}[1]{\cline{1-\value{spicols}} #1 \\}

% Options table: Name / Default / Range.
\newenvironment{optionstable}{%
  \begin{spitable}{3}{|p{80pt}|p{150pt}|p{140pt}|}%
    {\textbf{Name} & \textbf{Default} & \textbf{Range}}%
}{%
  \end{spitable}%
}
\newcommand{\optrow}[3]{\spirow{\texttt{#1} & #2 & #3}}

% Syntax block: verbatim, but \ { } still work as commands inside it.
% This needs the fancyvrb package; documents which only use tables do not
% need to load it, so the environment is only defined when it is available.
\ifdefined\Verbatim
\DefineVerbatimEnvironment{syntax}{Verbatim}{commandchars=\\\{\}}
\fi
\newcommand{\syntaxhead}{\par\textbf{SYNTAX:}\nopagebreak}
\newcommand{\syntaxdesc}{\emph{/** description */}}

% A variable name such as $returnType.
\newcommand{\spival}[1]{\texttt{\$#1}}

% Marks the macros as loaded, so generated files can load this file on demand:
%   \ifdefined\spidocloaded\else% spidoc.tex - macros used by the generated reference section.
% Load from the preamble with \input{spidoc}.
% Needs the longtable and fancyvrb packages, which the main document loads.

% General table.
%   #1 number of columns, #2 column specification, #3 heading row
%
% The table has one extra empty column on the right, and the rules are drawn
% with \cline over the real columns only. Rows written by hand may therefore
% end with "& \\" as they always did. (Once those rows have been tidied up the
% extra column can be dropped and \cline replaced by \hline.)
\newcounter{spicols}
\newenvironment{spitable}[3]{%
  \setcounter{spicols}{#1}%
  \small
  \begin{longtable}{#2l}
  \cline{1-#1} #3 \\ \endfirsthead
  \cline{1-#1} #3 \\ \cline{1-#1} \endhead
  \cline{1-#1} \multicolumn{#1}{r}{\emph{Continued on next page}} \\ \endfoot
  \cline{1-#1} \endlastfoot
}{%
  \end{longtable}%
}

% One row of a general table: the cells separated by &.
\newcommand{\spirow}[1]{\cline{1-\value{spicols}} #1 \\}

% Options table: Name / Default / Range.
\newenvironment{optionstable}{%
  \begin{spitable}{3}{|p{80pt}|p{150pt}|p{140pt}|}%
    {\textbf{Name} & \textbf{Default} & \textbf{Range}}%
}{%
  \end{spitable}%
}
\newcommand{\optrow}[3]{\spirow{\texttt{#1} & #2 & #3}}

% Syntax block: verbatim, but \ { } still work as commands inside it.
% This needs the fancyvrb package; documents which only use tables do not
% need to load it, so the environment is only defined when it is available.
\ifdefined\Verbatim
\DefineVerbatimEnvironment{syntax}{Verbatim}{commandchars=\\\{\}}
\fi
\newcommand{\syntaxhead}{\par\textbf{SYNTAX:}\nopagebreak}
\newcommand{\syntaxdesc}{\emph{/** description */}}

% A variable name such as $returnType.
\newcommand{\spival}[1]{\texttt{\$#1}}

% Marks the macros as loaded, so generated files can load this file on demand:
%   \ifdefined\spidocloaded\else\input{spidoc}\fi
\newcommand{\spidocloaded}{}

\fi
\newcommand{\spidocloaded}{}

\fi
\newcommand{\spidocloaded}{}

\fi
\newcommand{\spidocloaded}{}

